\chapter{ПОСТАНОВКА ЗАДАЧИ И ОПИСАНИЕ ИСХОДНЫХ ДАННЫХ} \label{ch1}

% не рекомендуется использовать отдельную section <<введение>> после лета 2020 года
%\section{Введение. Сложносоставное название первого параграфа первой главы для~демонстрации переноса слов в содержании} \label{ch1:intro}

В данной главе формулируется задача прогнозирования фенотипических 
признаков сельскохозяйственных растений по генотипическим данным, 
представленным в виде однонуклеотидных полиморфизмов (SNP). 
Приводится математическая постановка задачи многомерной регрессии, 
описывается обобщённая архитектура предлагаемой модели, выделяются 
её ключевые модификации по сравнению с известными аналогами. 
Кроме того, даётся характеристика трёх наборов генетических и 
фенотипических данных (нут, рис, рапс), используемых в вычислительных 
экспериментах. В \ref{ch1:sec1} представлена формальная постановка 
задачи, в \ref{ch1:sec2} приведено описание исходных данных.


\section{Постановка задачи}
\label{ch1:sec1}

Пусть имеется выборка из \(n\) образцов сельскохозяйственных растений. 
Каждый образец характеризуется генотипическими и фенотипическими данными.\cite{Lu2015}

\subsection{Генотипические данные}

SNP (ОНП) — это единичное изменение нуклеотида
в последовательности ДНК, которое встречается в популяции с частотой
не менее 1\%. SNP-маркеры являются наиболее распространённым типом
генетической изменчивости и используются в качестве молекулярных
маркеров для ассоциативных исследований и геномного прогнозирования.

Генотип \(i\)-го образца представляется вектором SNP-маркеров:
\[
\mathbf{x}_i = (x_{i1}, x_{i2}, \dots, x_{im})^\top \in \mathcal{X},
\]
где \(m\) — общее количество SNP. Каждый SNP кодируется числом
минорных аллелей:
\[
x_{ij} \in \{0, 1, 2\},
\]
где \(0\) соответствует мажорной гомозиготе (\(AA\)), \(1\) — гетерозиготе
(\(Aa\)), \(2\) — минорной гомозиготе (\(aa\)). Таким образом, значение
\(x_{ij}\) отражает дозу минорного аллеля в генотипе.

\subsection{Фенотипические данные}
Каждый образец имеет вектор фенотипических признаков:
\[
\mathbf{y}_i = (y_{i1}, y_{i2}, \dots, y_{ik})^\top \in \mathbb{R}^k,
\]
где \(k\) — количество прогнозируемых признаков (в данной работе 
рассматриваются случаи \(k = 2\) и \(k = 3\)).

\subsection{Задача многомерной регрессии}

Пусть задана обучающая выборка \(\{\mathbf{x}_i, \mathbf{y}_i\}_{i=1}^n\), где
\(\mathbf{x}_i \in \mathcal{X} \subseteq \mathbb{R}^m\) — вектор генотипов
(SNP-маркеров) \(i\)-го образца, \(\mathbf{y}_i \in \mathbb{R}^k\) — вектор
соответствующих фенотипов. Требуется построить модель
\[
f_\theta: \mathcal{X} \to \mathbb{R}^k,
\]
которая по вектору SNP-маркеров \(\mathbf{x}\) предсказывает вектор
фенотипов \(\hat{\mathbf{y}} = f_\theta(\mathbf{x})\). Параметры
\(\theta\) модели выбираются из условия минимизации эмпирического риска:
\[
\mathcal{L}(\theta) = \frac{1}{n} \sum_{i=1}^{n} \ell\bigl(f_\theta(\mathbf{x}_i),\; \mathbf{y}_i\bigr) \to \min_{\theta},
\]
где \(\ell(\cdot, \cdot)\) — функция потерь (в данной работе используется
среднеквадратическая ошибка MSE). Усреднение по \(n\) образцам позволяет
оценить математическое ожидание ошибки на генеральной совокупности и
обеспечивает независимость значения функции потерь от объёма выборки.

\subsection{Обобщённая архитектура модели}

Предлагаемая модель \(f_\theta\) состоит из четырёх последовательных
блоков, при этом итоговое предсказание формируется как сумма двух
компонент:
\[
f_\theta(\mathbf{x}) = P\bigl(T\bigl(E(\mathbf{x})\bigr)\bigr) \;+\; \lambda \cdot R\bigl(T\bigl(E(\mathbf{x})\bigr)\bigr),
\]
где:
\begin{itemize}
	\item \(E(\cdot)\) — блок генетического кодирования, преобразующий
	SNP-вектор \(\mathbf{x}\) в эмбеддинги генов;
	\item \(T(\cdot)\) — трансформерный блок для моделирования межгенных
	взаимодействий;
	\item \(P(\cdot)\) — блок прогнозирования фенотипов (основная задача);
	\item \(R(\cdot)\) — реконструкционный блок (автокодировщик),
	выступающий в роли регуляризатора.
\end{itemize}

Символ «\(+\)» в выражении для \(f_\theta(\mathbf{x})\) означает
покомпонентное сложение двух векторов одинаковой размерности: вектора
фенотипов из блока \(P\) и вектора реконструированных признаков из
блока \(R\), умноженного на коэффициент \(\lambda\). В контексте
многозадачного обучения такое сложение позволяет:
\begin{enumerate}
	\item одновременно оптимизировать основную задачу (предсказание
	фенотипов) и вспомогательную задачу (реконструкция входных SNP);
	\item использовать реконструкционный блок как дополнительный
	регуляризатор, способствующий обучению более информативных
	эмбеддингов генов;
	\item с помощью коэффициента \(\lambda\) регулировать вклад
	регуляризации в общую функцию потерь.
\end{enumerate}

В данной работе реконструкционный блок используется только на этапе
обучения, а при инференсе (прогнозировании) его выход отключается
(\(\lambda = 0\)).

\subsection{Ключевые модификации}
По сравнению с базовой архитектурой в данной модели реализованы 
следующие усовершенствования:
\begin{enumerate}
	\item \textbf{Механизм внимания на уровне SNP} — адаптивное 
	взвешивание маркеров внутри каждого гена;
	\item \textbf{Реконструкционная регуляризация} — дополнительная 
	потеря на восстановление исходных SNP-данных;
	\item \textbf{Многозадачная функция потерь с адаптивным 
		взвешиванием} — автоматическая балансировка вклада различных 
	фенотипов через обучаемые параметры неопределённости.
\end{enumerate}

\subsection{Анализ существующих решений}

Современные методы геномного прогнозирования можно разделить на две основные группы: классические статистические подходы и методы машинного, в частности глубокого, обучения.

\paragraph{Статистические методы.}
Традиционно для предсказания фенотипов по SNP-данным используются линейные модели, такие как rrBLUP (ridge regression best linear unbiased prediction) \cite{Endelman2011}, GBLUP (genomic best linear unbiased prediction) и байесовские подходы (BGLR) \cite{Perez2014}. Эти методы предполагают, что эффекты отдельных SNP аддитивны и независимы, что позволяет эффективно работать с данными, где количество маркеров значительно превышает число наблюдений. Однако они имеют ограниченную способность улавливать нелинейные и эпистатические взаимодействия между генами \cite{Crossa2017}. Кроме того, методы отбора признаков, такие как анализ главных компонент, могут использоваться для снижения размерности, но при этом теряют часть биологически значимой информации \cite{Manthena2022}.

\paragraph{Методы глубокого обучения.}
В последние годы всё большее распространение получают нейросетевые архитектуры, способные автоматически извлекать сложные иерархические признаки из геномных данных. Свёрточные нейронные сети (CNN), как в работе \cite{Ji2024}, эффективны для выявления локальных паттернов, но не всегда учитывают глобальные взаимосвязи между удалёнными генетическими регионами. Результаты исследований показывают, что CNN могут быть успешно применены для прогнозирования признаков нута с использованием комбинации SNP и физических расстояний \cite{Bankin2024}.

\paragraph{Трансформерные архитектуры.}
Особый интерес представляют модели на основе трансформеров, которые благодаря механизму многоголового внимания \cite{Vaswani2017} позволяют моделировать сложные взаимодействия между любыми элементами входной последовательности. В задаче геномного прогнозирования трансформеры продемонстрировали высокую эффективность для предсказания устойчивости ячменя к фузариозу \cite{Jubair2021} и для выявления зависимостей между SNP-маркерами \cite{Wu2023}. Обобщение этих подходов рассмотрено в работе \cite{Consens2023}, где обсуждается применение больших языковых моделей для анализа генома.

\paragraph{Работа Реутова и базовые архитектуры.}
Непосредственным предшественником данной работы является исследование \cite{Reutov2025}, в котором была разработана модель прогнозирования признаков нута (масса тысячи семян, высота растения, биомасса) с использованием трансформерной архитектуры. В этой работе была предложена агрегация SNP-маркеров по генам с помощью свёртки, а также исследованы методы аугментации данных (shuffle и mixup). Сравнительный анализ показал, что метод mixup с долей аугментированных данных 0.4 позволяет повысить точность предсказания фенотипов, особенно для сложных признаков. Однако в модели \cite{Reutov2025} использовалась свёрточная агрегация SNP, не учитывающая относительную важность отдельных маркеров внутри гена. Это ограничение послужило одним из мотивов для разработки предлагаемой в данной работе архитектуры с механизмом внимания на уровне SNP.

\paragraph{Выявление генных взаимодействий.}
Для интерпретации моделей глубокого обучения и выявления значимых генов и их взаимодействий широко применяются значения Шепли \cite{Shapley1953, Lundberg2017}. В задаче детекции парных генных взаимодействий этот подход был адаптирован и показал свою эффективность для выявления синергетических и антагонистических эффектов \cite{Cui2022}. В работе предшественника \cite{Reutov2025} значения Шепли были использованы для оценки центральности генов и построения матриц взаимодействий для фенотипов TSW и Ptht нута, что позволило идентифицировать ключевые гены-«хабы».

\paragraph{Обоснование предлагаемых модификаций.}
Таким образом, анализ существующих решений выявляет следующие нерешённые проблемы и направления для улучшения:

\begin{enumerate}
	\item Большинство трансформерных моделей для геномного прогнозирования используют агрегацию SNP-маркеров на уровне генов без учёта индивидуальной важности каждого маркера (свёртка или усреднение).
	\item Регуляризация скрытых представлений генов остаётся недостаточно изученной; стандартные методы (dropout, weight decay) не всегда предотвращают переобучение при малых выборках.
	\item Отсутствуют систематические сравнения стратегий отбора генов (взвешенной, равномерной по геному, равномерной по хромосомам) для трансформерных архитектур.
	\item Интерпретация выявленных генных взаимодействий требует дальнейшей валидации и сопоставления с биологическими базами данных.
\end{enumerate}

Предлагаемая в данной работе модель направлена на решение этих проблем за счёт введения механизма внимания на уровне SNP, реконструкционной регуляризации и адаптивного взвешивания фенотипических признаков.

\section{Описание данных}
\label{ch1:sec2}
\subsection{Данные по нуту}

В исследовании использовались генотипические и фенотипические данные по образцам нута (\textit{Cicer arietinum}), собранные в различные годы: 2016, 2017 и 2022 (два набора --- Астраханская опытная станция (AOS) и Кубанская опытная станция (KOS)). Данные были предоставлены Всероссийским институтом растениеводства им. Н.И. Вавилова (ВИР) и обработаны в Санкт-Петербургском политехническом университете Петра Великого (СПбПУ). Подробное описание этих данных приведено в работе \cite{Reutov2025}, где они использовались для первоначальной апробации базовой архитектуры.

Каждый образец характеризуется набором из 5783 однонуклеотидных полиморфизмов (SNP), аннотированных к конкретным генным регионам, а также тремя количественными фенотипическими признаками:
\begin{itemize}
	\item \textbf{TSW (Thousand Seed Weight)} --- масса тысячи семян (г);
	\item \textbf{Ptht (Plant Height)} --- высота растения (см);
	\item \textbf{Bvld (Biological Yield)} --- биомасса всего растения (г).
\end{itemize}

SNP-маркеры закодированы численно в соответствии с типом аллеля: 2 --- минорная гомозигота ($aa$), 1 --- гетерозигота ($Aa$), 0 --- мажорная гомозигота ($AA$). Данные по фенотипам Ptht и Bvld представлены средними значениями по пяти растениям для каждого образца.

Для обучения и тестирования модели использовались данные 2016 и 2017 годов как наиболее согласованные по распределению признаков. Обучающая выборка составила 442 образца, тестовая --- 99 образцов. Статистические характеристики данных приведены в таблице \ref{tab:chickpea_stats}.

%\begin{table}
%	\centering
%	\caption{Сводная статистика по датасетам нута}
%	\label{tab:chickpea_stats}
%	\begin{tabular}{|l|c|c|c|c|}
%		\hline
%		\textbf{Датасет} & \textbf{Образцов} & \textbf{TSW (среднее ± std)} & \textbf{Ptht (среднее ± std)} & \textbf{Диапазон TSW} \\ \hline
%		2016    & 407 & 223.35 ± 86.84 & 54.34 ± 8.31 & [90.0, 475.0] \\ \hline
%		2017    & 146 & 222.09 ± 83.11 & 53.66 ± 6.12 & [60.0, 435.0] \\ \hline
%		2022\_AOS & 113 & 245.98 ± 84.94 & 50.67 ± 9.99 & [69.8, 450.6] \\ \hline
%		\textbf{Всего} & \textbf{666} & \textbf{226.92 ± 85.96} & \textbf{53.57 ± 8.30} & \textbf{[60.0, 475.0]} \\ \hline
%	\end{tabular}
%\end{table}

\begin{table}
	\centering
	\small
	\caption{Сводная статистика по датасетам нута}
	\label{tab:chickpea_stats}
	\begin{tabular}{|l|c|c|c|c|}
		\hline
		\textbf{Датасет} & \textbf{Образцов} & \textbf{TSW (среднее ± std)} & \textbf{Ptht (среднее ± std)} & \textbf{Диапазон TSW} \\ \hline
		2016    & 407 & 223.35 ± 86.84 & 54.34 ± 8.31 & [90.0, 475.0] \\ \hline
		2017    & 146 & 222.09 ± 83.11 & 53.66 ± 6.12 & [60.0, 435.0] \\ \hline
		2022\_AOS & 113 & 245.98 ± 84.94 & 50.67 ± 9.99 & [69.8, 450.6] \\ \hline
		\textbf{Всего} & \textbf{666} & \textbf{226.92 ± 85.96} & \textbf{53.57 ± 8.30} & \textbf{[60.0, 475.0]} \\ \hline
	\end{tabular}
	\normalsize
\end{table}

\subsection{Данные по рису}

Для оценки обобщающей способности разработанной архитектуры были использованы генотипические и фенотипические данные риса (\textit{Oryza sativa}). Данные получены из открытых источников \cite{Lu2015} и предварительно обработаны для обеспечения совместимости с архитектурой модели.


Датасет по рису содержит следующие компоненты:

\begin{enumerate}
	\item \textbf{Фенотипические данные}: 
	\begin{itemize}
		\item Количество образцов: 469
		\item Фенотипические признаки:
		\begin{itemize}
			\item Высота растения (\textit{Plant height}) [см]
			\item Число побегов (\textit{Tiller number})
		\end{itemize}
		\item Дополнительные признаки: угол побега, параметры флагового листа, параметры метёлки и др.
	\end{itemize}
	
	\item \textbf{Генотипические данные}:
	\begin{itemize}
		\item Количество образцов: 523
		\item Количество SNP-маркеров: 4\,045
		\item Формат кодирования: численное представление аллелей (0, 1, 2)
		\item Идентификаторы SNP: \texttt{gwseq-rs1}, \texttt{gwseq-rs10}, \texttt{seq-rs992} и др.
	\end{itemize}
	
	\item \textbf{Аннотационные данные}:
	\begin{itemize}
		\item Координаты SNP: 4\,136 записей с информацией о хромосоме и позиции
		\item Генные позиции: 211\,705 генов с указанием хромосомы, начала и конца
		\item Сопоставление SNP с генами: 15\,653 соответствий между SNP и генами
	\end{itemize}
\end{enumerate}

\paragraph{Статистические характеристики данных}

После объединения генотипических и фенотипических данных и фильтрации общих образцов получен датасет размерностью $469 \times 4\,045$ SNP-признаков и 2 целевых фенотипических признаков.

Статистические характеристики фенотипических признаков представлены в таблице \ref{tab:rice_stats}.

\begin{table}[h]
	\centering
	\caption{Статистические характеристики фенотипических данных риса}
	\label{tab:rice_stats}
	\begin{tabular}{|l|c|c|}
		\hline
		\textbf{Параметр} & \textbf{Plant height (cm)} & \textbf{Tiller number} \\ \hline
		Количество образцов & 469 & 469 \\ \hline
		Минимальное значение & 54.50 & 4.67 \\ \hline
		Максимальное значение & 199.50 & 26.50 \\ \hline
		Среднее значение & 110.93 & 12.43 \\ \hline
		Стандартное отклонение & 24.57 & 3.18 \\ \hline
		Медиана & 111.20 & 12.00 \\ \hline
		Коэффициент вариации, \% & 22.15 & 25.61 \\ \hline
	\end{tabular}
\end{table}

Для высоты растения наблюдается распределение, близкое к нормальному, с центром в районе 110 см. Распределение числа побегов имеет положительную асимметрию с концентрацией значений в диапазоне 10--15 побегов.

\paragraph{Корреляционная структура данных}

Между целевыми фенотипическими признаками наблюдается умеренная отрицательная корреляция. Коэффициент корреляции Пирсона составил:
\[
r_{\text{Plant height} \times \text{Tiller number}} = -0.348 \quad (p < 0.001)
\]
Это указывает на существование биологической связи между высотой растения и числом побегов: более высокие растения имеют тенденцию к образованию меньшего числа побегов.

\paragraph{Предобработка данных}

Перед подачей данных в модель необходимо выполнить ряд шагов предобработки, которые являются неотъемлемой частью разработанной архитектуры.

\subparagraph{Частотное кодирование SNP}

Значения SNP, изначально закодированные как \(\{0, 1, 2\}\), преобразуются в частоту встречаемости соответствующего аллеля в обучающей выборке. Для каждого SNP-маркера \(j\) вычисляется:
\[
f_{j}(a) = \frac{n_{j}(a)}{N_{\text{train}}}, \quad a \in \{0, 1, 2\},
\]
где \(n_{j}(a)\) — количество образцов в обучающей выборке со значением \(a\) для SNP \(j\), \(N_{\text{train}}\) — общее количество образцов в обучающей выборке. Тогда преобразованное значение для образца \(i\) и SNP \(j\) составляет:
\[
x'_{ij} = f_{j}(x_{ij}) = \frac{n_{j}(x_{ij})}{N_{\text{train}}}.
\]
Такое кодирование позволяет модели учитывать редкость или распространённость конкретного аллеля в популяции.

\subparagraph{Объединение данных и разделение на выборки}

После кодирования выполняется объединение генотипических и фенотипических данных по идентификаторам образцов. Далее данные разделяются на три выборки: 80\% — обучающая, 10\% — валидационная, 10\% — тестовая.

\subparagraph{Аугментация методом MixUp}

Для увеличения объёма обучающей выборки и улучшения обобщающей способности модели применяется метод MixUp. Для случайно выбранной пары образцов \(( \mathbf{x}_i, \mathbf{y}_i )\) и \(( \mathbf{x}_j, \mathbf{y}_j )\) синтетический пример генерируется следующим образом:
\[
\begin{aligned}
	\tilde{\mathbf{x}} &= \lambda \mathbf{x}_i + (1 - \lambda) \mathbf{x}_j, \\
	\tilde{\mathbf{y}} &= \lambda \mathbf{y}_i + (1 - \lambda) \mathbf{y}_j,
\end{aligned}
\]
где \(\lambda \in [0, 1]\) — параметр смешивания. В данной работе используется фиксированное значение \(\lambda = 0.0625\).

\subparagraph{Преобразование SNP $\rightarrow$ Gene}

Ключевым этапом предобработки, специфичным для разработанной архитектуры, является преобразование данных с уровня отдельных SNP на уровень генов. Для этого на основе таблицы аннотации SNP-маркеров к генам каждому SNP сопоставляется соответствующий ген. В результате формируется отображение SNP $\rightarrow$ Gene, которое используется блоком генетического кодирования для агрегации SNP-признаков в эмбеддинги генов.

Данное преобразование позволяет:
\begin{itemize}
	\item сократить размерность признакового пространства;
	\item учесть биологическую иерархию генетической информации;
	\item обеспечить интерпретируемость модели на уровне генов, а не отдельных SNP.
\end{itemize}

Для данных по рису результаты преобразования выглядят следующим образом:
\begin{itemize}
	\item Всего сопоставлений SNP-ген: 15\,653;
	\item Уникальных генов: 3\,918;
	\item Среднее количество SNP на ген: $\approx 4.0$;
\end{itemize}
Таким образом, размерность признакового пространства сокращается с 4\,045 SNP-признаков до 3\,918 ген-признаков. Важно отметить, что плотность SNP на ген в данных риса составляет всего около 1.07 после фильтрации аннотированных маркеров, что создаёт дополнительные сложности для агрегации.

\subparagraph{Итоговые размерности данных}

После выполнения всех шагов предобработки получены следующие размерности:
\begin{itemize}
	\item Признаки: \(\mathbf{X} \in \mathbb{R}^{469 \times 4045}\)
	\item Целевые переменные: \(\mathbf{y} \in \mathbb{R}^{469 \times 2}\)
\end{itemize}
\subsection{Данные по рапсу}

Для исследования эффективности предложенных архитектурных решений на данных с высокой плотностью SNP-маркеров были использованы генотипические и фенотипические данные рапса (\textit{Brassica napus}). Данные получены из базы CropGS-Hub \cite{Chen2024} (кохорта GSTP010), содержащей результаты полногеномного секвенирования 991 образца из 39 стран.


SNP-маркеры в данных CropGS-Hub имеют имена, отражающие их локализацию. Например:
\begin{center}
	\texttt{chrA01\_17073}
\end{center}
где \texttt{chrA01} — хромосома (A01 рапса), а \texttt{17073} — позиция на хромосоме в парах оснований.

Основные характеристики генетических данных:
\begin{itemize}
	\item Количество образцов: 991
	\item Количество SNP: 1\,272\,066 уникальных
	\item Количество генов с аннотацией: 60\,768
	\item Покрытие аннотацией: 100\%
	\item Среднее количество SNP на ген: 21.3
	\item Медиана SNP на ген: 16
	\item Максимум SNP на ген: 745
\end{itemize}

Данные соответствуют описанию в статье \cite{Wu2019}, где 991 образец рапса были секвенированы со средней глубиной 6.6×, а SNP вызывались с использованием эталонного генома Darmor-bzh.

\paragraph{Распределение SNP по генам}

Распределение количества SNP на ген (рис. \ref{fig:snp_per_gene_rapeseed}) является асимметричным: большинство генов содержат от 3 до 28 SNP, при этом встречаются гены с очень большим числом маркеров (до 745). Такая вариативность критически важна для блока кодирования, где механизм внимания должен адаптивно обрабатывать гены с разным количеством SNP.

\begin{figure}
	\centering
	\includegraphics[width=0.9\textwidth]{snp_per_gene_analysis.png}
	\caption{Распределение количества SNP на ген в данных рапса}
	\label{fig:snp_per_gene_rapeseed}
\end{figure}

\paragraph{Распределение по хромосомам}

Данные охватывают 19 хромосом (A01--A10, C01--C09). Распределение SNP и генов по хромосомам представлено на рисунке \ref{fig:chromosome_distribution_rapeseed}. Наибольшее количество SNP сосредоточено на хромосомах A03 и C03, что соответствует их длине и плотности генов. Плотность SNP на ген варьируется от 21.7 до 24.0, что говорит о примерно равномерном покрытии генов маркерами по хромосомам.

\begin{figure}
	\centering
	\includegraphics[width=0.9\textwidth]{chromosome_distribution.png}
	\caption{Распределение SNP и генов по хромосомам рапса}
	\label{fig:chromosome_distribution_rapeseed}
\end{figure}

\paragraph{Позиции SNP внутри генов}

Для позиционного кодирования важны относительные позиции SNP внутри генов. Как показано на рисунке \ref{fig:relative_position_rapeseed}, распределение близко к равномерному (среднее 0.5003, медиана 0.5007), что является благоприятным фактором для использования тригонометрических функций в позиционном кодировании.

\begin{figure}
	\centering
	\includegraphics[width=0.8\textwidth]{relative_position_analysis.png}
	\caption{Распределение относительных позиций SNP внутри генов рапса}
	\label{fig:relative_position_rapeseed}
\end{figure}

\paragraph{Расстояния до генов}

Значительная часть SNP (60\,236 SNP) находятся в промоторной области (расстояние до старта гена > 0). На рисунке \ref{fig:distance_analysis_rapeseed} представлено распределение расстояний до начала и конца гена. Медианное расстояние до старта составляет 1482 п.н., что может быть учтено при построении позиционного кодирования как дополнительный признак.

\begin{figure}
	\centering
	\includegraphics[width=0.8\textwidth]{distance_analysis.png}
	\caption{Распределение расстояний до начала и конца гена (логарифмическая шкала)}
	\label{fig:distance_analysis_rapeseed}
\end{figure}

\paragraph{SNP, относящиеся к нескольким генам}

В данных обнаружено 22\,706 SNP (1.78\% от общего числа), которые отнесены к нескольким генам. Это явление может быть обусловлено перекрывающимися генами, расположением SNP в регуляторных областях, особенностями аннотации генома, а также аллополиплоидной природой рапса (наличие A и C субгеномов). В данной работе для обработки таких SNP используется стратегия выбора гена с минимальным расстоянием до SNP.

\paragraph{Фенотипические признаки}

Фенотипические данные рапса включают широкий спектр признаков (табл. \ref{tab:rapeseed_traits}):

\begin{table}
	\centering
	\caption{Основные фенотипические признаки рапса в данных CropGS-Hub}
	\label{tab:rapeseed_traits}
	\begin{tabular}{|l|l|}
		\hline
		\textbf{Категория} & \textbf{Признаки} \\ \hline
		Время цветения & FloweringTime (от всходов до цветения) \\ \hline
		Качество семян & OilContent, ProteinContent, GlucosinolateContent \\ \hline
		Жирные кислоты & C16:0, C18:1, C18:2, C18:3, C20:1, C22:1 \\ \hline
		Другие & Высота растения, масса семян и др. \\ \hline
	\end{tabular}
\end{table}

Для данной работы наибольший интерес представляют признаки, сопоставимые с исследованными ранее для нута: время цветения (FloweringTime) коррелирует с фенологическими сроками, содержание масла (OilContent) и белка (ProteinContent) отражают продуктивность и качество семян, а высота растения (Plant height) является важным морфологическим признаком.


\paragraph{Пропуски в данных}

На рисунке \ref{fig:missing_values_rapeseed} показана доля пропущенных значений по признакам. Большинство признаков имеют пропуски на уровне 70–80\%, что связано с тем, что не для всех линий были проведены соответствующие биохимические анализы. Исключение составляет время цветения (\texttt{FloweringTime}), для которого пропуски составляют лишь 6.6\% (65 из 991 образца). Это делает данный признак наиболее пригодным для обучения и послужило основанием для использования его в качестве основной целевой переменной.

\begin{figure}
	\centering
	\includegraphics[width=0.8\textwidth]{rapeseed_missing_values.png}
	\caption{Доля пропущенных значений по признакам рапса}
	\label{fig:missing_values_rapeseed}
\end{figure}



\subsection{Сравнительная характеристика наборов данных}

Для наглядного сопоставления трёх использованных наборов данных в таблице \ref{tab:datasets_comparison} сведены их основные параметры.

\begin{table}
	\centering
	\small
	\caption{Сравнительная характеристика наборов данных}
	\label{tab:datasets_comparison}
	\begin{tabular}{|l|c|c|c|}
		\hline
		\textbf{Параметр} & \textbf{Нут} & \textbf{Рис} & \textbf{Рапс} \\ \hline
		Образцов & 553 & 469 & 991 \\ \hline
		SNP (после фильтрации) & 5\,783 & 1\,708 & 1\,272\,066 \\ \hline
		Генов (уникальных) & 2\,837 & 1\,599 & 60\,768 \\ \hline
		Среднее SNP/ген & 2.04 & 1.07 & 21.3 \\ \hline
		Покрытие аннотацией & 100\% & 42\% & 100\% \\ \hline
		Основные признаки & \begin{tabular}{@{}c@{}}TSW, Ptht, Bvld\end{tabular} & \begin{tabular}{@{}c@{}}Plant height,\\ Tiller number\end{tabular} & \begin{tabular}{@{}c@{}}FloweringTime,\\ OilContent,\\ GlucosinolateContent\end{tabular} \\ \hline
		Источник & ВИР / СПбПУ & \begin{tabular}{@{}c@{}}Lu et al. (2015)\end{tabular} & \begin{tabular}{@{}c@{}}CropGS-Hub\\ (Wu et al., 2019)\end{tabular} \\ \hline
	\end{tabular}
	\normalsize
\end{table}

Как видно из таблицы, наборы данных существенно различаются по плотности SNP на ген: от крайне низкой у риса (1.07) до высокой у рапса (21.3). Это позволяет провести кросс-культурный анализ и оценить, как плотность маркеров влияет на эффективность предложенных архитектурных решений, в частности механизма внимания на уровне SNP.

\subsection{Выводы по главе \thechapter}

В данной главе сформулирована задача многомерной регрессии для прогнозирования фенотипов по SNP-данным. Представлена обобщённая архитектура модели, включающая четыре блока: кодирования с вниманием на уровне SNP, трансформерный, прогнозирования и реконструкционный. Описаны три набора данных (нут, рис, рапс), существенно различающиеся по плотности SNP на ген и набору целевых признаков, что позволит в дальнейшем провести кросс-культурный анализ и оценить обобщающую способность разработанной модели.


%% Вспомогательные команды - Additional commands
%
%\newpage % принудительное начало с новой страницы, использовать только в конце раздела
%\clearpage % осуществляется пакетом <<placeins>> в пределах секций
%\newpage\leavevmode\thispagestyle{empty}\newpage % 100 % начало новой страницы